% Generated draft: requires mathematical and source review.
\documentclass[11pt]{article}
\usepackage[margin=25mm]{geometry}
\usepackage{fontspec}
\setmainfont{Latin Modern Roman}
\usepackage{amsmath,amssymb,mathtools,unicode-math}
\setmathfont{Latin Modern Math}
\usepackage{enumitem,needspace,xurl,hyperref}
\setlength{\parindent}{0pt}
\setlength{\parskip}{5pt}
\setlength{\emergencystretch}{3em}
\newcommand{\N}{\mathbb{N}}
\newcommand{\Z}{\mathbb{Z}}
\newcommand{\Q}{\mathbb{Q}}
\newcommand{\R}{\mathbb{R}}
\newcommand{\C}{\mathbb{C}}
\newcommand{\F}{\mathbb{F}}
\newcommand{\A}{\mathbb{A}}
\newcommand{\PP}{\mathbb{P}}
\newcommand{\E}{\mathbb{E}}
\newcommand{\eps}{\varepsilon}
\newcommand{\dd}{\,\mathrm{d}}
\begin{document}
\section*{UnsolvedMath ID: 30004315}
\textbf{Source problem code:} catalogue entry
\subsection{Definitions and mathematical statement}
For a set \(X\), a set-theoretic solution of the Yang--Baxter equation, in the braid convention used here for exposition, is a map \(r\colon X\times X\to X\times X\) satisfying
\[
(r\times\operatorname{id}_X)\circ(\operatorname{id}_X\times r)\circ(r\times\operatorname{id}_X)
=
(\operatorname{id}_X\times r)\circ(r\times\operatorname{id}_X)\circ(\operatorname{id}_X\times r)
\]
as maps \(X^3\to X^3\). This definition imposes no bijectivity or nondegeneracy requirement; whether such additional requirements belong to the intended problem is not specified in the supplied record.

A skew lattice is an algebra \((L,\wedge,\vee)\) whose two binary operations are associative and idempotent and satisfy, for all \(x,y\in L\),
\[
 x\wedge(x\vee y)=x=(y\vee x)\wedge x,
 \qquad
 x\vee(x\wedge y)=x=(y\wedge x)\vee x.
\]
Commutativity of the operations is not required. The supplied record does not define the broader term ``noncommutative lattice.''

\textbf{Problem as supplied.} Do the Yang--Baxter solutions obtained from skew lattices remain solutions for other noncommutative lattices, and do noncommutative lattices yield new solutions?

The first question asks whether the constructions referred to in the source continue to satisfy the Yang--Baxter equation when applied to suitable other noncommutative lattice structures. The second asks whether such structures give additional solutions. A fully quantified mathematical formulation remains unresolved: the record supplies neither the solution formulas nor their original hypotheses, the axioms of the proposed broader classes, or the meaning of ``new.'' Consequently, no universal assertion over arbitrary noncommutative lattices, and no specific construction or novelty criterion, is asserted here.
\subsection{Short English statement}
The problem asks whether constructions of Yang--Baxter solutions from skew lattices extend to other noncommutative lattice structures and whether these structures produce new solutions. The supplied record leaves the constructions, broader classes, additional solution requirements, and meaning of ``new'' unspecified.
\subsection{Sources}
\begin{itemize}
\item Supplied local UnsolvedMath registry record, permanent ID \(30004315\), title ``Yang--Baxter Solutions from Noncommutative Lattices,'' source problem code \texttt{catalogue entry}, in the Oberwolfach Reports Open Problems collection. The record gives proposed year 2019 and names no individual proposer.
\item Source report, as cited in the supplied record: \emph{Mini-Workshop: Algebraic Tools for Solving the Yang--Baxter Equation} (2020), Oberwolfach Reports, pp. 3207--3242, DOI 10.4171/owr/2019/51. \url{https://doi.org/10.4171/owr/2019/51}. The supplied background describes these reports as documenting mathematical workshops and their open-problem discussions.
\item Supplied literature reference: \emph{Skew lattices and set-theoretic solutions of Yang--Baxter}. \url{https://doi.org/10.1016/j.jalgebra.2019.10.007}.
\item Supplied literature reference: \emph{Generalized digroups and di-skew braces}. \url{https://link.springer.com/article/10.1007/s00233-025-10606-2}.
\item Supplied literature reference: \emph{Generalized semi-braces and Yang--Baxter solutions}. \url{https://arxiv.org/abs/2004.01606}.
\item Supplied literature assessment, checked on \texttt{2026-08-22}: the record labels the problem open and states that the skew-lattice case is established, that broader brace-like structures yield additional solutions, and that no general theorem covers arbitrary noncommutative lattices. This is supplied dated curation derived from the OWR material, not an independently verified literature assessment. The record explicitly calls for independent verification of time-sensitive status and cited claims.
\end{itemize}
\end{document}
